\section{Partial configurations and their deterministic estimates}
\label{sec:foundations}

Two kinds of estimates are needed before pairing can proceed. The
structural accounts control actual paths and separate the unresolved
stubs. The analytic estimates turn trace and entry bounds for the
current inverse into the quadratic and moment bounds needed for one
transition. We first derive the common matrix and path identities, then
develop these two sets of estimates. The analytic implications are
conditional: Sections~\ref{sec:early-phase} and~\ref{sec:cleanup} establish
their hypotheses inductively along the selected pairing history.

Throughout this section, an operator on $\mathbf 1^\perp$ is extended by
zero when vertex entries are taken. All constants are uniform in $n$;
they may depend on the fixed parameters chosen in the parameter section.
The notation $n^{o(1)}$ below represents explicitly computable products of
fixed powers of $\log n$ and $\exp(O((\log\log n)^2))$.

\subsection{The configuration pencil and a single exposure}

We begin by assigning a matrix to each partial pairing. Its dependence
on the unresolved stubs will make the change caused by one exposure
explicit.

Give every vertex $d$ distinguished stubs. For a partial pairing, let $S$
be its $l$ unresolved stubs. For $l>0$, let $C_S$ be the vertex--stub
incidence matrix, and put
\[
 t=C_S\mathbf 1,\qquad
 \Pi_l=I-\frac{\mathbf 1\mathbf 1^T}{l},\qquad
 \Xi=C_S\Pi_l,\qquad \mathcal D=\Xi\Xi^T.
\]
Write $\xi_x$ for the column of $\Xi$ indexed by $x$. If $A$ is the
adjacency matrix of the resolved pairs, define
\[
 A_e=A+\frac{tt^T}{l},\qquad
 M(h)=z(h)I-\sigma A_e-h\mathcal D,\qquad
 m_h=\frac{h}{1-h^2},\qquad W=M(h)^{-1}/m_h.
\]
The operators $M(h)$ and $W$ act on $\mathbf 1^\perp$, and $W$ is used
only when $M(h)$ is positive definite. At the empty pairing, $W$ is the
identity on $\mathbf1^\perp$. At a complete pairing, define
$M(h)=z(h)I-\sigma A_G$ and $W=M(h)^{-1}/m_h$ on this subspace;
positivity for both signs
therefore bounds the nonconstant spectrum in absolute value by $z(h)$.
The partial-configuration pencil lets us maintain this positivity during
the exposure. We have $A_e\mathbf 1=d\mathbf 1$
and, writing $\alpha=1-qh^2$,
\[
 -\partial_hM(h)=\mathcal D+h^{-2}\alpha I.
\]

The next lemma identifies the inverse with a formal path generating
function and computes its update when one pair is resolved. Both
descriptions will be used: the path expansion controls entries, and the
matrix update controls the change in the inverse.
For the relation between nonbacktracking walks and the quadratic
Ihara--Bass matrix polynomial, see
\cite[Section~6, Eq.~(6.2)]{FZ1999}. The recurrence below
also includes unresolved averaging blocks, and we prove it directly.

\begin{lemma}[Path generating function and exposure identities]
\label{lemma:configuration-algebra}
Let $C$ be the incidence matrix of all $dn$ stubs. On resolved stubs let
$P_*$ exchange partners, and on $S$ let it be the averaging projection.
Set $D_e=CP_*^2C^T=dI-\mathcal D$ and
\[
 Q_0=I,\qquad
 Q_j=CP_*[(C^TC-I)P_*]^{j-1}C^T\quad(j\ge1).
\]
Then, as a formal power series on $\mathbf 1^\perp$,
\[
 W=\sum_{j\ge0}(\sigma h)^jQ_j.
\]
For distinct $u,v\in S$, $l>2$, define
\[
 U=[\xi_u,\xi_v],\quad
 G_0=I-J_2/l,\quad B=G_0^{-1}=I+J_2/(l-2),\quad
 H_0=-\sigma h FB,
\]
where $J_2$ is the all-ones matrix and $F$ exchanges the two coordinates.
If this pair is resolved, then
\[
 \mathcal D'=\mathcal D-UBU^T,\qquad
 A_e'=A_e+UFBU^T,\qquad
 M'=M+U(hI-\sigma F)BU^T.
\]
With $G=U^TWU$, $D_p=G-G_0$, and
$H=(H_0^{-1}+D_p)^{-1}$, the inverse update, whenever its inverses
exist, is
\[
 W'=W-WUHU^TW.
\]
Uniformly for the prescribed compact range of $h$, if $\|D_p\|=o(1)$,
then $M'>0$ whenever $M>0$.
\end{lemma}

\begin{proof}
We first prove the generating-function identity, then the exposure
formulas and the assertion about positivity.

The matrix $Q_j$ counts weighted nonbacktracking stub walks. A resolved
traversal is deterministic; an unresolved traversal lands uniformly in
$S$, independently on each use. After landing, the arrival stub is
excluded from the next departure. Thus every row sum is $dq^{j-1}$ for
$j\ge1$. Since $P_*^3=P_*$,
\[
 CP_*^2(C^TC-I)P_*=(D_e-I)CP_*.
\]
Expansion gives $Q_1=A_e$, $Q_2=A_e^2-D_e$, and
$Q_j=A_eQ_{j-1}-(D_e-I)Q_{j-2}$ for $j\ge3$. Consequently
\[
 \sum_{j\ge0}(\sigma h)^jQ_j
 =(1-h^2)[I-\sigma hA_e+h^2(D_e-I)]^{-1}=W.
\]
For the exposure identities, write $a=e_{v(u)}$, $b=e_{v(v)}$ and
replace $t$ by $t-a-b$ and $l$ by $l-2$ in
$\mathcal D=\operatorname{diag}(t)-tt^T/l$ and
$A_e=A+tt^T/l$. Direct expansion gives the displayed formulas.
The Woodbury inverse formula \cite{HS1981} applies because
\[
 [m_h(hI-\sigma F)B]^{-1}+G_0=H_0^{-1}.
\]
It remains to verify preservation of positivity.
Both $m_h(hI-\sigma F)B$ and $H_0^{-1}$ have one positive and one
negative eigenvalue, bounded away from zero uniformly in the prescribed
range. The perturbation $D_p=o(1)$ preserves the latter inertia.

We recall the block congruence behind the Schur-complement argument.
For real symmetric blocks $\mathsf A,\mathsf D$ with $\mathsf A$
invertible, and a compatible real rectangular block $\mathsf B$, define
\[
 \mathsf L=
 \begin{bmatrix}I&-\mathsf A^{-1}\mathsf B\\0&I\end{bmatrix}.
\]
Direct multiplication gives
\[
 \mathsf L^T
 \begin{bmatrix}\mathsf A&\mathsf B\\\mathsf B^T&\mathsf D\end{bmatrix}
 \mathsf L
 =\begin{bmatrix}
 \mathsf A&0\\0&\mathsf D-\mathsf B^T\mathsf A^{-1}\mathsf B
 \end{bmatrix}.
\]
Since $\mathsf L$ is invertible, the inertia of the original block
matrix is the sum of the inertias of $\mathsf A$ and
$\mathsf D-\mathsf B^T\mathsf A^{-1}\mathsf B$; here inertia records
the numbers of positive, negative, and zero eigenvalues. When
$\mathsf D$ is invertible, the same identity with the block order
exchanged gives the other Schur complement.

With
$K=(hI-\sigma F)B$, apply the two Schur complements to
$\left[\begin{smallmatrix}M&U\\U^T&-K^{-1}\end{smallmatrix}\right]$.
One complement is $M'$ and the other is
$-m_h(H_0^{-1}+D_p)$, proving that the inertia of $M'$ equals that of
$M$.
\end{proof}

Let $P=P_{<L}$ be the degree-$<L$ truncation of this generating function,
and write $E=W-P$. The formal exposure subtracts the degree-$<L$
truncation of
\begin{equation}
 P U\left[H_0\sum_{j\ge0}
       (-(U^TPU-G_0)H_0)^j\right]U^TP.
 \label{eq:original-1}
\end{equation}
This statement is coefficientwise: every factor $H_0$ has positive
degree, so every coefficient uses only finitely many terms, and every
coefficient of degree less than $L$ is unchanged on replacing $W$ by
$P$.

We will repeatedly use a cancellation that is exact, including the
finite-population correction. For two scalar arrays $a,c$ on $S$, take
an ordered distinct pair uniformly in $I\subseteq S$, where $|I|=l-k$.
Write $a_p=(a(u),a(v))^T$ and $c_p=(c(u),c(v))^T$ for its two-entry
arrays. Their contraction satisfies
\begin{equation}
 \begin{aligned}
 \frac{\mathbb E[a_p^TH_0c_p]}{-2\sigma h}
 &=\frac{(1+1/(l-2))(\sum_Ia)(\sum_Ic)}{|I|(|I|-1)}\\
 &\quad-\frac{k\sum_Iac}{(l-2)|I|(|I|-1)}.
 \end{aligned}
 \label{eq:original-2}
\end{equation}
Indeed, the swap part has average
$2[(\sum_Ia)(\sum_Ic)-\sum_Iac]/[|I|(|I|-1)]$, while the diagonal part
contributes $2\sum_Iac/|I|$. Substituting the coefficients of $FB$
gives \eqref{eq:original-2}. In particular the expectation vanishes for
centered arrays when $I=S$.

\subsection{Structural accounts and separated unresolved stubs}

The path bounds will require control of small graph cores. We construct
accounts for these cores and for proximity among unresolved stubs. The
proximity and return counts identify the stubs that cleanup must resolve
to obtain separation. We verify that cleanup preserves the core bounds.

In the early phase, expose a uniform ordered distinct pair whose
vertices have actual graph distance greater than $g_0$. The graph on
stubs of rejected pairs has maximum degree
$D_g=O_d(q^{g_0})$. Every resulting actual graph is simple and has girth
at least $g_0$.

\begin{lemma}[Structural accounts]
\label{lemma:structure}
Fix $J=\lfloor .9\log_qn\rfloor$ and a sufficiently large fixed integer
$r_*$. The joint discounted-account rule can maintain, through the
early and cleanup phases, the following bound for every connected
multigraph core $H$ of minimum degree at least two, excess
$r=e_H-v_H+1\in\{1,\ldots,r_*\}$, and $s_H=e_H\le r_*J$:
\begin{equation}
 \#\{\text{stub embeddings of }H\}
       \le n^{1-r+o(1)}q^{s_H}.
 \label{eq:original-3}
\end{equation}
An embedding is injective on vertices and uses distinct proper stubs
for all incidences at a vertex. The accounts and their next-step
averages are computable in polynomial time.

At $l=l_1$, fewer than $b$ unresolved stubs are bad, where a stub is bad
if its vertex has a nonempty actual nonbacktracking return of length at
most $R_0$, or lies within distance $R_0$ of another unresolved stub
(including distance zero). Cleanup can resolve all these bad stubs and
stop at $l_0$ with one unresolved stub per unresolved vertex, and no
nonempty actual nonbacktracking path of length at most $R_0$ joining
two unresolved vertices, including a return to one vertex.
\end{lemma}

\begin{proof}
We first construct the core accounts and bound their change under
rejection in the early phase. These bounds give a global path count and,
in particular, the short-return count needed to identify bad stubs. We
then bound proximity, perform cleanup, and verify that cleanup preserves
the same core and path estimates.

\emph{Initial core accounts.}
Suppressing degree-two vertices in a core leaves a bounded-size
multigraph, apart from the separate case of a cycle. Its subdivisions
with at most $r_*J$ edges give $(1+\log n)^{O(r_*)}$ shapes. Cores with
a vertex of degree exceeding $d$ need not be included.

We use the following form of the perfect-hash-family theorem
\cite[Section~4.4, Theorem~3(iii)]{NSS1995}: an $n$-element set has an explicitly
constructible
$k$-perfect family of maps to $k$ colors, of size
$\exp(k) k^{O(\log k)}\log n$, constructible in
$\exp(k) k^{O(\log k)}n^{O(1)}$ time. This is the specialization
$\ell=k$ of their splitter theorem: splitting a $k$-element set
equally among $k$ colors makes the map injective on that set.
The bounded case $k=1$ uses a single constant map.
For $k=v_H$, sum the unrestricted
completion probabilities over embeddings colorful for each map.
Every embedding is included at least once. If an embedding is already
incompatible with the partial pairing it contributes zero; if it has
$p$ demanded pairs still pending, it contributes
\[
 f_p(l)=\prod_{i=1}^p(l-2i+1)^{-1}.
\]
For one coloring the number of colorful ordered vertex assignments is
at most $k!(n/k)^k$. At the empty pairing the proper stub assignments
contribute $\prod_x(d)_{\deg_H(x)}$, where $(d)_j$ is a falling
factorial. Since $e_H\ge v_H$ and $(d)_j\le dq^{j-1}$,
\[
 \frac{\prod_x(d)_{\deg_H(x)}}{d^{e_H}}
 \le d^{v_H-e_H}q^{2e_H-v_H}\le q^{e_H}.
\]
Stirling's estimate cancels the factor $\exp(k)$ against $k!/k^k$,
leaving only a polynomial factor in $k$. Since $k=O_d(\log n)$,
the remaining factor $k^{O(\log k)}\log n$ is $n^{o(1)}$.
The initial account is therefore at most
$n^{1-r+o(1)}q^{s_H}$.

We use color-coding on a bounded-width tree decomposition, following
Alon, Yuster and Zwick \cite[Section~6, Theorem~6.3]{AYZ1995}; the states below also retain
prescribed stub assignments and pending demanded-pair counts. These
colorful sums do not require enumeration of
$n^k$ assignments. Choose a spanning tree of $H$ and a width-one
tree decomposition of that tree. Add the endpoints of the $r$
non-tree edges to every bag. Every edge of $H$ is then covered,
the bags containing each vertex remain connected, and the width
is at most $2r+1=O(r_*+1)$. In a bag record its vertex and stub
assignments, the subset of colors used in the processed part, and the
number of pending demanded pairs. At a join the color subsets must be
disjoint off the bag. At an edge check its two stub statuses and
increment the pending count when appropriate. After totaling by that
count, multiply by $f_p(l)$. There are
$2^{O(k)}n^{O(r_*+1)}$ states and transitions, including joins; all
counts and probabilities have polynomial bit length.

\emph{Change under rejection.}
Under unrestricted exposure this account is a martingale. It remains
to bound the cost of restricting to accepted pairs. Conditional on an
eventual matching
containing $p$ demanded pairs, the next exposed pair is one of these
pairs with total probability $2p/l$; otherwise it is a uniform pair
on the other $l-2p$ stubs. Deleting $2p$ vertices from a rejected-pair
graph of maximum degree $D_g$ changes its rejected-pair proportion by
$O(pD_g/l^2)$. The contribution of demanded hits to a possible
increase in acceptance probability is also $O(pD_g/l^2)$. Bayes'
formula thus bounds the account's conditional expectation by
$1+O(pD_g/l^2)$ times its present value. Summing to $l_1$ costs
$\exp(O(r_*JD_g/l_1))=1+o(1)$. Normalize by
$n^{1-r}q^{s_H}$ and apply the joint account rule.

\emph{Global path diagrams and short returns.}
At every early-phase state, \eqref{eq:original-3} with $r_*>10$,
increased if necessary, bounds a core of excess $r_*$ and at most
$r_*J$ edges by $n^{1-.1r_*+o(1)}<1$. Thus no such actual core occurs.

A connected union of a fixed number of actual nonbacktracking walks,
with total length at most $r_*J$ and a
fixed number of marked endpoints, has excess less than $r_*$. Otherwise
retain a connected subgraph of excess $r_*$ and prune its leaves,
contradicting the absence of such cores just proved. Suppressing unmarked
degree-two vertices leaves boundedly many vertices, so the number of
diagram shapes is $n^{o(1)}$. A diagram with $e'$ edges and excess $r$
has at most
\[
 n^{1-r+o(1)}q^{e'}
\]
embeddings: apply \eqref{eq:original-3} to the core and extend its
bounded-leaf trees at cost $Cq^{e'-s_H}$; for a tree use $Cnq^{e'}$.
The number of walk descriptions in an embedded diagram is bounded by
a constant depending on the fixed parameters. Indeed, branching is
possible only at boundedly many marked or branch vertices, and two
visits to the same such vertex are separated by at least $g_0$ steps.
The ratio $r_*J/g_0$ is bounded independently of $n$.

In particular, the union traced by a nonempty nonbacktracking return
of length $k$ contains a cycle, so its excess is at least one and it has
at most $k$ edges. The global diagram count therefore gives at most
$n^{o(1)}q^k$ such return walks. Summing over $k\le R_0$ and allowing
at most $d$ unresolved stubs at each vertex gives at most
$n^{o(1)}q^{R_0}$ return-bad stubs. This deduction uses only the
early-phase core bound and girth exclusion already established.

\emph{Proximity at the end of the early phase.}
For proximity at $l_1$, count simple paths of lengths $0$ through
$R_0$ with two distinguished, distinct endpoint stubs required to
survive until $l_1$. For an assignment with $p$ pending required pairs
and $r$ marked surviving stubs, its completion probability is
\[
 \frac{(l_1)_r\,2^p((l-l_1)/2)_p}{(l)_{r+2p}},
\]
with value zero for incompatibility. This follows by assigning stubs
to the future pair slots and the final surviving slots of a uniform
permutation. Multiply the contribution by $w^p$, with
$w=1+CD_g/l_1$. A demanded next hit removes one factor $w$ and hence
pays its entire rejection inflation. Conditional on a successful
non-hit, the pair is uniform on the ordinary stubs; deleting the
$2p+r$ distinguished stubs costs only
$1+O((p+r)D_g/l^2)$. The initial normalized colorful sum is
$n^{o(1)}$: without normalization its expectation is
\[
 n^{o(1)}O\left(q^{R_0}\operatorname{polylog}(n)l_1^2/n\right).
\]
For positive-length paths this follows by choosing proper incidences,
paying $(dn)^{-j}(1+o(1))$ for the $j$ edges and
$O((l_1/n)^2)$ for the two surviving flags; length zero obeys the same
bound. Together with the return-bad count above, this bounds all bad stubs.
Both contributions are $o(b)$, since
\[
 q^{R_0}l_1^2/n=n^{1/3+.34a+o(1)},\qquad
 b=n^{1/3+.345a+o(1)}.
\]

\emph{Cleanup.}
The preceding bounds leave enough steps to resolve every bad stub.
Fix the bad set at $l_1$. Resolve each bad stub as first endpoint,
choosing its partner uniformly among surviving initially clean stubs.
Then use any surviving first stub and choose its partner uniformly
among the eligible surviving initially clean stubs until $l_0$; there
are exactly $b$ steps available. An untouched clean
stub is initially farther than $R_0$ from every other unresolved stub.
A path of length at most $R_0$ starting there cannot reach a new
cleanup edge: before its first such edge it would already reach an
initial unresolved vertex in the initial graph. Its short-path and
short-return properties therefore persist. Every cleanup pair is
admissible, and the asserted final separation follows.

It remains to preserve the core accounts during these cleanup steps.
Multiply $f_p(l)$ by
$(1+Cb/l_0)^p$. If the fixed first stub is demanded, success requires
its unique demanded partner; deletion of at most $O(b)$ possible
partners is paid by the lost weight. Otherwise a successful exposure
has probability ratio at most
$f_p(l-2)/f_p(l)=1+O(p/l)$, irrespective of the omissions. The total
discount cost is $\exp(O(r_*Jb/l_0))=1+o(1)$, as is the change on
reinitializing the weights. This completes the simultaneous account
construction. The core bound and girth exclusion therefore hold
through cleanup, so the global diagram and short-return estimates proved
above remain valid there as well.
\end{proof}

\subsection{Path diagrams, absolute majorants, and polynomial tails}

The coefficient majorants and their first-power row bound below hold
in the general partial-configuration setting. The stronger path estimates
are conditional on the current early or cleanup state having $l\ge l_0$,
the core bound \eqref{eq:original-3}, and actual girth at least $g_0$.
Lemma~\ref{lemma:structure} maintains these properties through both
phases; its final separation conclusion is used in
Section~\ref{sec:boundary-completion}. An admissible pair has endpoint
vertices at actual distance greater than $g_0$.

The global diagram bound proved with the structural accounts counts
actual paths over all choices of endpoints. For entrywise inverse
updates we also need bounds with endpoints fixed, and absolute
majorants that include virtual traversals and centering. These
majorants will then control the tail discarded from the formal exposure.

We begin with fixed endpoints. The union of all walks
of length at most $L$ has bounded excess. After the first walk, every walk adding an edge also
increases excess: an added forest would have a new leaf, whereas the
only possible leaves are the already present endpoints. At most
$r_*+1$ walks would witness excess $r_*$. Since $L/J$ is bounded below
one, taking $r_*$ sufficiently large ensures
$(r_*+1)L\le r_*J$, a contradiction. The preceding branch-count
argument consequently bounds the number of walks at each length
between fixed endpoints by a constant. Likewise, a family of
cycle-containing diagrams rooted at $i$, with no possible leaf other
than $i$ and with each total length at most $J$, has bounded excess:
each edge-adding member increases it, including the first member.

\begin{lemma}[Absolute path bounds]
\label{lemma:path-majorants}
In the partial-configuration setting above, let an output label denote
either a vertex coordinate projected onto $\mathbf1^\perp$, or a current
centered stub $\xi_x$. For a label $i$ write
\[
 p_i(x)=P_{i,\xi_x},\quad e_i(x)=E_{i,\xi_x},\quad
 w_i(x)=W_{i,\xi_x}.
\]
We distinguish polynomial coefficient majorants from their numerical
evaluations. Let $N_k(u,v)$ count actual nonbacktracking stub walks of
length $k$ from $u$ to $v$, with $N_0(u,v)=\mathbf1_{u=v}$, and define
the following polynomials with nonnegative coefficients:
\[
 \begin{aligned}
 \mathsf U_l(z)&=\frac{C}{l}
       \left(1+\sum_{k=1}^{L-1}kq^kz^k\right),\\
 \mathsf A_{i,x}(z)&=\sum_{k=0}^{L-1}N_k(v(i),v(x))z^k
                         +\mathsf U_l(z),\\
 \mathsf J_x(z)&=\sum_{k=1}^{L-1}N_k(v(x),v(x))z^k
                         +\mathsf U_l(z),\qquad
 \mathsf B_{x,y}(z)=\mathsf A_{x,y}(z)+C/l.
 \end{aligned}
\]
Here $v(i)$ is the underlying vertex of the output label. The constant
$C$ is large enough to include the virtual, centering, and degree-zero
Gram corrections. Coefficientwise absolute values of $p_i(x)$ are
bounded by $\mathsf A_{i,x}$; diagonal and cross errors of $U^TPU-G_0$
are bounded by $\mathsf J_x$ and $\mathsf B_{x,y}$, respectively.
The actual-walk counts are symmetric, as are these majorants whenever
the input and output are stub labels. Write
\[
 A_i(x)=\mathsf A_{i,x}(h),\qquad
 j_x=\mathsf J_x(h),\qquad b_{xy}=\mathsf B_{x,y}(h).
\]
In particular $\mathsf U_l(h)\le C(1+L)^2q^{L/2}/l$, and
\begin{equation}
 \sum_{x\in S}A_i(x)\le C(1+L)^2q^{L/2}.
 \label{eq:path-row-sum}
\end{equation}

For the following bounds, assume that the current early or cleanup
state has $l\ge l_0$, satisfies \eqref{eq:original-3}, and has actual
girth at least $g_0$. Uniformly through $l_0$,
\begin{equation}
 \begin{split}
 &\sum_{x\in S}A_i(x)^2\le C,\\
 &\sum_{x\in S}j_x^2\le n^\eta,\qquad
   \sum_{x\in S}j_x\le n^\eta q^{L/2}.
 \end{split}
 \label{eq:original-4}
\end{equation}
For an admissible pair in either phase there is a fixed $c_0'>0$ with
\begin{equation}
 j_u+j_v+b_{uv}\le n^{-c_0'}.
 \label{eq:original-5}
\end{equation}
The same decay holds after multiplying nonconstant diagonal
coefficients by any fixed power of their length, including for global
diagonal entries with their degree-zero norm as baseline.
\end{lemma}

\begin{proof}
\emph{Coefficient majorants and the first row sum.}
An entry involving at least one virtual traversal at length $k$ is at
most $Ckq^k/l$: mark one such traversal, sum prefixes by the row sums,
pay its uniform landing factor $1/l$, and sum suffixes backwards using
reversibility and the same row sums. This overcounts and hence is an
upper bound. Vertex centering contributes $dq^{k-1}/n$; averaging an
endpoint against the unresolved stubs contributes at most $Ckq^k/l$.
The length-zero corrections are $O(1/l+1/n)$. These estimates prove the
coefficient majorants and the stated bound on $\mathsf U_l(h)$.
Reversing a walk proves symmetry. For \eqref{eq:path-row-sum}, each
vertex has at most $d$ unresolved stubs, so the actual length-$k$ row
sum is at most $C_dq^k$. Summing with weights $h^k$ and adding
$l\mathsf U_l(h)$ proves the claim. These arguments use no core or
girth hypothesis.

For the remaining estimates we use the stated structural hypotheses.

\emph{Actual-walk row squares.}
First consider actual walks from a fixed vertex to unresolved stubs.
In the squared weighted sum, identical walks contribute at most $d$:
choose the initial departure uniformly among $d$ stubs and successive
departures uniformly among $q$ alternatives, stopping at the first
unresolved departure. A stopped walk of $k$ actual steps has probability
$d^{-1}q^{-k}$, which dominates $d^{-1}h^{2k}$.

For two distinct walks ending at the same unresolved stub vertex,
remove their common terminal segment. Summing the possible common
continuations with squared weights again costs at most $d$ by the
stopping argument. The remaining pair forms a cycle-containing
diagram rooted at the initial vertex, with no other possible leaf.
If its total length $u$ is at most $J$, the rooted-family argument
above bounds its number by $n^{o(1)}$, and $u\ge g_0$ makes its total
weighted contribution $o(1)$. If $J<u\le2L$, take a sufficiently large
fixed $p$-th moment over roots. In a union of $p$ such diagrams every
edge-adding diagram increases excess, so a union of excess $r$ has at
most $\min(p,r)u$ edges. The global embedding bound gives a moment at
most
\[
 n^{1+o(1)}(1+q^u/n)^p.
\]
Here $r_*$ is chosen large enough to cover all $p$ diagrams. Thus each
root's count is at most $n^{1/p+o(1)}(1+q^u/n)$. After multiplication
by $q^{-u/2}$ this is bounded by
\[
 n^{1/p+o(1)}(q^{-J/2}+q^L/n)=o(1)
\]
for, for example, a sufficiently large fixed $p>4$. Summation over
the polylogarithmically many length choices preserves the decay. This proves
the actual row-square bound, with constant depending only on $d$ once
$n$ is large enough.

Adding the virtual and centering terms costs
$O((1+L)^4q^L/l)=o(1)$ in row-square sum, by $l\ge l_0$ and the
prescribed relation between $L$ and $l_0$. Thus the actual-walk bound
also proves the first estimate in \eqref{eq:original-4}.

\emph{Return bounds.}
The global number of return walks of length $k$ is at most
$n^{o(1)}q^k$, proving the first-power bound on $j$. For its square,
pair two return walks of lengths $k,k'$. Excess at least two costs at
most $n^{-1+o(1)}q^{(k+k')/2}$ after weighting, whose sum is $o(1)$.
At excess one the union consists of a root-to-cycle tail and the
cycle. Each return covers all these edges, so their number is at most
$\min(k,k')$; its weighted embedding bound is $n^{o(1)}$. This proves
the return-square bound, after absorbing subpolynomial factors in the
prescribed $n^\eta$.

\emph{Decay at admissible pairs.}
For an admissible pair, girth and endpoint distance exclude all short
nonconstant terms. Bounded fixed-endpoint multiplicity bounds the
remaining sum by $C\sum_{k\ge g_0}q^{-k/2}$, in addition to the
power-small uniform errors. The degree-zero centered Gram error is
$O(1/l)$, even with other stubs at either vertex. Fixed powers of $k$
do not change the resulting positive decay exponent.
\end{proof}

\paragraph{The formal exposure tail.}
The coefficientwise majorants also control the part discarded when
the formal exposure is truncated. At a current early or cleanup state
with $l\ge l_0$, the core bound \eqref{eq:original-3}, and actual girth
at least $g_0$, let $u,v$ be any distinct unresolved stubs. The absolute
coefficient tail of \eqref{eq:original-1}, of total degree at least $L$,
has every entry on global coordinates or current centered-stub labels
bounded by
\begin{equation}
 B_{\rm tail}:=Cq^{-(1/2-\epsilon)L}
       \le n^{.30a}/\sqrt{l_0},
 \label{eq:original-6}
\end{equation}
where $\epsilon>0$ is fixed and sufficiently small, and $C$ may depend
on $\epsilon$.

To prove this, evaluate the nonnegative coefficient majorants at
$z_h=hq^{1/2-\epsilon}\le q^{-\epsilon}<1$. Put
$r=\lfloor g_0/3\rfloor$. Between two fixed vertices there is at most
one actual nonbacktracking walk of length less than $r$, counting all
lengths together: each such walk is simple, and two distinct ones would
form a cycle shorter than $g_0$. This includes the degree-zero walk
when the vertices coincide. Its contribution is at most one. Bounded
fixed-endpoint multiplicity bounds the longer contribution by
$Cz_h^r/(1-z_h)=o(1)$. The rescaled uniform error also tends to zero:
\[
 \mathsf U_l(z_h)\le C(1+L)^2q^{(1-\epsilon)L}/l=o(1),
 \qquad q^L/l_0=n^{-.58a+o(1)}.
\]
Consequently all outer majorants are at most $1+o(1)$. A nonnegative
coefficient majorant $\mathsf D(z_h)$ for the two-by-two middle
deviation $U^TPU-G_0$ has diagonal entries $o(1)$ and off-diagonal
entries at most $1+o(1)$. The degree-zero Gram corrections are $O(1/l)$;
the possible degree-zero off-diagonal contribution for two stubs at the
same vertex is included in the bound by one. Since
\[
 \|\mathsf D(z_h)\|_\infty\le1+o(1),\qquad
 \bigl\lVert z_h(F+J_2/(l-2))\bigr\rVert_\infty
 =z_h\left(1+\frac2{l-2}\right),
\]
the product of the middle and $H_0$ majorants has norm at most
$q^{-\epsilon}(1+o(1))$. For sufficiently large $n$ this is bounded
by a fixed $\rho_\epsilon<1$. Coefficientwise nonnegative matrix
multiplication therefore bounds the whole formal series at $z_h$ by
$C_\epsilon$. Returning to $h$ pays
$q^{-(1/2-\epsilon)L}$ for every term of degree at least $L$.
Finally $q^{-L/2}\asymp n^{.29a}/\sqrt{l_0}$, so choosing
$\epsilon$ sufficiently small proves \eqref{eq:original-6}.
Every fixed power of the number of factor positions is also summable,
since $\sum_{m\ge0}(m+1)^k\rho_\epsilon^m<\infty$ for fixed $k$.
This is a bound on the formal $P$ exposure; it does not assert positivity
after an arbitrary pair is resolved. The actual transitions use
admissible pairs, for which \eqref{eq:original-5} and the entry caps give
the stronger power-small ratio in their numerical mixed expansions.

\subsection{A coefficient bound from positivity of the pencil}

The path estimates control entries of the truncation. Positivity at
both signs supplies an operator bound for each coefficient and hence
for the whole truncation. We derive that bound from the scalar
polynomials obtained by evaluating the pencil on unit vectors.

\begin{lemma}[Polynomial operator bounds]
\label{lemma:pencil-coefficients}
If the precision is positive for both signs at $h$, then
\begin{equation}
 \|h^jQ_j|_{\mathbf1^\perp}\|\le C(1+j)^2,
 \qquad \|P\|\le C(1+L)^3.
 \label{eq:original-7}
\end{equation}
For a fixed stub endpoint, the signed column of $P$ over global
vertices has polylogarithmic squared norm.
For the following absolute column and entry bounds, also assume that
the current early or cleanup state has $l\ge l_0$, satisfies the core
bound \eqref{eq:original-3}, and has actual girth at least $g_0$.
For every fixed integer $p\ge3$, the absolute $p$-power sum of that
column is bounded by a constant; for
sufficiently large $n$ that constant can depend only on $d$ and $p$.
Stub columns have bounded absolute squared sums, and every outer
entry of $P$ in the formal exposure expansion \eqref{eq:original-1} has absolute
value at most $1+o(1)$.
\end{lemma}

\begin{proof}
\emph{Operator bounds.}
On $\mathbf1^\perp$ put
$F(z)=I-zA_e+z^2(qI-\mathcal D)$. For a complex unit vector $v$,
$v^*F(z)v=1-xz+yz^2$ has real $x,y$ and $y\le q$. A nonreal pair of
roots has modulus $y^{-1/2}\ge q^{-1/2}$. A real root in $(-h,h)$ is
also impossible. Endpoint positivity would require another root on
the same side of zero in that interval, but their product would be
less than $h^2\le1/q$, contradicting $1/y\ge1/q$. Linear polynomials
and polynomials with roots of opposite signs are excluded directly
by the endpoint signs. Thus every root has modulus at least $h$.
Factoring the scalar polynomial shows
\[
 |v^*F(z)v|\ge(j+1)^{-2}
 \quad\text{when }|z|=h\left(1-\frac1{j+1}\right),\quad j\ge1.
\]
Since $\|F(z)v\|\ge|v^*F(z)v|$, its inverse has norm at most
$(j+1)^2$. Apply the Cauchy coefficient formula to
$(1-z^2)F(z)^{-1}$; the extra radius ratio is at most a fixed
constant. This proves \eqref{eq:original-7}. Moreover,
$\|\xi_x\|\le C$, so
\[
 \sum_i|P_{i,\xi_x}|^2=\|P\xi_x\|^2
 \le C\|P\|^2\le C(1+L)^6,
\]
where $i$ ranges over global vertices. This proves the signed
column-square bound from positivity alone.

\emph{Column and entry bounds.}
We now combine the operator estimate with the path bounds.
For actual paths of length $j$ to a fixed endpoint, bounded
multiplicity and the total path count give an $\ell^3$ norm
$O(q^{-j/6})$ after weighting, which is summable. Uniform errors have
total cube sum
$O(n[(1+L)^2q^{L/2}/l]^3)=o(1)$. This proves the cubic bound. For the
refined constant, paths shorter than $g_0/3$ are unique in total at
fixed endpoints; their cube sums are at most
$1+\sum_{j\ge1}dq^{j-1}q^{-3j/2}$. Longer actual paths have power-small
supremum and absolute aggregate squared sum $O((1+L)^2)$, so their
$p$-power sums tend to zero. The uniform correction has vanishing
$p$-power sum by its cube bound and vanishing supremum. Apply the triangle inequality in
$\ell^p$. The same short/long separation proves the $1+o(1)$ entry
bound. The cubic column bound has a constant depending only on $d$.
Combining it with the entry bound gives, for every fixed integer $p\ge3$
and all sufficiently large $n$ (depending on $p$),
\[
 \sum_i |P_{i,\xi_x}|^p
 \le (1+o(1))^{p-3}\sum_i |P_{i,\xi_x}|^3
 \le 2C_d.
\]
Here $i$ ranges over global vertex coordinates. Thus the large-power
column constants needed in the early transition
argument can be fixed before its scalar moment order. The same
short/long separation also shows that choosing a fixed distance radius
from a fixed positive entry threshold does not require fixing the later
structural losses first.
\end{proof}

\subsection{Deterministic consequences of the resolvent bounds}
\label{subsec:resolvent-consequences}

This subsection derives deterministic consequences of the trace and entry
bounds in \eqref{eq:original-8}. The transition arguments in the following
sections establish those bounds by induction. At each transition these
implications are applied only to the current state. We first bound the
norm and spectral tail of a positive weighted inverse. We then use
nearby parameter values to obtain local quadratic and averaged trace
estimates, from which the required moments follow. Set
$F_1=\operatorname{Tr}(\mathcal DW)$ and $F_0=\operatorname{Tr}W$.
We use the trace notation
\[
 \begin{gathered}
 \Gamma=\Xi^TW\Xi,\qquad L_1=F_1-(l-1),\qquad L_0=F_0-(n-1),\\
 K_{\rm all}=\Gamma^2,\qquad Q_{\rm all}=\Xi^TW^2\Xi,\qquad
 \mathcal T=\operatorname{Tr}Q_{\rm all}.
 \end{gathered}
\]
At the primaries, the squared-resolvent trace $\mathcal T$ will be kept
away from zero by applying the account rule to its reciprocal.
At all tracked primary and auxiliary points suppose positivity and
\begin{equation}
 |F_1-(l-1)|\le n^{.09a}/\alpha_*,\qquad
 |F_0-(n-1)|\le n^{.12a}/\alpha_*^2,\qquad
 \max_{i,j}|E_{ij}|\le n^{.37a}/\sqrt{l_0},
 \label{eq:original-8}
\end{equation}
where the entries range over the specified global and centered-stub
labels. At a current early or cleanup state with $l\ge l_0$, the core
bound \eqref{eq:original-3}, and actual girth at least $g_0$,
Lemma~\ref{lemma:pencil-coefficients} bounds the formal $P$ outer entries
by $1+o(1)$. The corresponding numerical $W$ outer entries have the
same bound: $W=P+E$, and the entry cap in \eqref{eq:original-8} is
$n^{.37a}/\sqrt{l_0}=o(1)$. Write $v_*=.125a$ and $c=.03a$.

\begin{lemma}[Consequences of the resolvent bounds]
\label{lemma:local-laws-implications}
At a current early or cleanup state, assume $l\ge l_0$, the core
bound \eqref{eq:original-3}, and actual girth at least $g_0$.
Assume positivity and \eqref{eq:original-8} at every tracked primary
and auxiliary parameter, for both signs.
For $B_h^+=\mathcal D+h^{-2}\alpha I$, the positive operator
$S_h=(B_h^+)^{1/2}W(B_h^+)^{1/2}$ has norm at most
$n^{v_*}/\alpha$. For a sufficiently large fixed $C$, the sum of its
eigenvalues above $Cn^c/\alpha$ is at most $n^{v_*}/\alpha$.
These bounds hold at the primaries and their auxiliaries. At all these
points, every global coordinate or centered stub $f$ satisfies
\begin{equation}
 f^TW(\mathcal D+h^{-2}\alpha I)Wf
   \le(1+h^{-2})\|f_\perp\|^2+n^{-c_1}.
 \label{eq:original-10}
\end{equation}
At the primaries, the averaged trace identity is
\begin{equation}
 \frac{\operatorname{Tr}K_{\rm all}}l+
       h^{-2}\alpha\frac{\mathcal T}{l}
       =1+h^{-2}+O(n^{-c_2}).
 \label{eq:original-11}
\end{equation}
The following moment bounds hold at both primaries and auxiliaries:
\begin{equation}
 \begin{split}
 \operatorname{Tr}K_{\rm all}^2+
       \alpha^2\operatorname{Tr}Q_{\rm all}^2
   &\le C\left(n^{2c}l/\alpha^2+n^{4v_*}/\alpha^4\right),\\
 \operatorname{Tr}(\Gamma K_{\rm all})+
       \alpha\operatorname{Tr}(\Gamma Q_{\rm all})
   &\le C\left(n^cl/\alpha+n^{3v_*}/\alpha^3\right),\\
 \|Q_{\rm all}\|&\le Cn^{2v_*}/\alpha^3.
 \end{split}
 \label{eq:original-12}
\end{equation}
The decay exponents $c_1,c_2>0$ are fixed independently of $n$.
\end{lemma}

\begin{proof}
\emph{Reduction to the closest primary.}
We begin with the norm and spectral tail. Positivity at the closest
primary controls the other primary groups directly; the closest group
will require the auxiliary trace bounds.

For every primary group except the closest one, positivity at the
closest point and integration of $-\partial_hM(h)$ give
\[
 M(h)\succeq c_d\alpha(\mathcal D+\alpha I).
\]
Indeed the interval between these points has length comparable with
$\alpha$, and on a fixed subinterval its value of $1-qt^2$ is
comparable with $\alpha$. The claim follows there with norm $C/\alpha$
and an empty tail. It remains to control the closest point.

\emph{A positive finite-difference filter at the closest primary.}
At that point put $R_h=M(h)^{-1}=m_hW$,
$X=(B_h^+)^{-1/2}M(h)(B_h^+)^{-1/2}$ and $S=X^{-1}$. Positivity,
\eqref{eq:original-8}, and $\operatorname{Tr}W=O(n)$ give a polynomial
upper bound on $S$. The bounds $B_h^+\succeq c_d\alpha_*I$ and
$\|M(h)\|=O(1)$ give an inverse-polynomial lower bound. Set
$k=k_h$ and $b=b_h=n^{-c}\alpha_*$.

For the linear pencil $X+xI$, the alternating trace difference at
$0,b,\ldots,kb$ is the nonnegative quantity
\begin{equation}
 \mathcal B:=\sum_{\lambda\in\operatorname{spec}S}
   \lambda\prod_{j=1}^k\frac{jb\lambda}{1+jb\lambda}.
 \label{eq:original-9}
\end{equation}
For $b\lambda$ small, the product suppresses the contribution of
$\lambda$; for $\lambda\ge C/b$, it is bounded below by a positive
constant depending only on $k,C$. A bound on $\mathcal B$ therefore
bounds the spectral tail that we need. We will obtain that bound from
the trace laws for the actual pencil, by comparing its alternating
trace difference with this linear one.

Write the actual pencil in the form $X+xI+\gamma(x)H_*$, where
\[
 H_*=(B_h^+)^{-1},\quad
 \gamma(x)=z(h-x)-z(h)-h^{-2}\alpha_*x,\quad
 T(x)=\operatorname{Tr}[X+xI+\gamma(x)H_*]^{-1}.
\]
Taylor's formula gives $\gamma(x)=O(x^2)$,
$\gamma'(x)=O(x)$, and bounded higher fixed derivatives on the
interval used. At $x=jb_h$, $0\le j\le k_h$, the trace laws imply
\[
 T(x)=m_{h-x}[(l-1)+h^{-2}\alpha_*(n-1)]
          +O(n^{.12a}/\alpha_*).
\]
Consequently,
\[
 \left|\sum_{j=0}^k(-1)^j\binom kjT(jb)\right|
 \le C_k(n^{.12a}/\alpha_*+nb^k).
\]

\emph{Comparison with the linear pencil.}
We show that the preceding trace difference is
$\mathcal B+o(1)+o(\mathcal B)$. The matrices $X$ and $H_*$ need not
commute, so the comparison retains the order of all matrix factors. For
$X_g(x)=X+xI+\gamma(x)H_*$ and $R(x)=X_g(x)^{-1}$,
the divided-difference product rule applied to $RX_g=I$ gives
\[
 R[x_0,\ldots,x_k]
 =\sum_{0=i_0<\cdots<i_p=k}(-1)^p
 R(x_0)\prod_{t=1}^p
 \left(X_g[x_{i_{t-1}},\ldots,x_{i_t}]R(x_{i_t})\right),
\]
where products retain their displayed order and $x_j=jb$. At nonzero
nodes expand about $A_j=(X+jbI)^{-1}$. Because
$\|\gamma(jb)H_*A_j\|\le Cb/\alpha_*=Cn^{-c}$, every such Neumann
series can be truncated at a fixed order $K$ with residual norm
at most $C_Kb^{-1}(Cb/\alpha_*)^K$. After multiplication by $b^k$,
every remaining factor, rank, and initial inverse norm is bounded by
a fixed power of $n$. Choosing $K$ sufficiently large makes the sum
of all residual traces $o(1)$.

After multiplication by $(-1)^kk!b^k$ and taking the trace, the term
with all first divided differences equal to $I$ and all nodes
unperturbed is exactly $\mathcal B$. It remains to show that all other
terms have total absolute value $o(1)+o(\mathcal B)$.
Use a common decomposition into the $O(\log n)$ dyadic spectral
blocks of $S$ for $S$ and every $A_j$, all of which are functions
of $S$. In a nonzero block tuple let
$\lambda_0$ be its largest scale, and let $r_0$ be the rank of a
block at that scale. A cyclic trace is bounded by $r_0$ times the
product of operator norms, by starting the trace at that block.
This choice is available even when the largest block occurs in an
intermediate inverse: cyclicity moves that block to the beginning of
the trace while retaining the cyclic order of all factors.
Write $s_0=\min(\lambda_0,b^{-1})$. The corresponding main scale is,
within constants depending on $k$, $r_0\lambda_0(bs_0)^k$.
This is an upper comparison scale for a mixed block tuple, rather
than the value of that tuple. The initial inverse block has norm at
most $C\lambda_0$, and every nonzero-node inverse block has norm
at most $C_k s_0$, irrespective of where the largest block occurs.
The contribution to $\mathcal B$ from that largest spectral block
is comparable to this scale, since on its dyadic interval
$\lambda\prod_{j=1}^k jb\lambda/(1+jb\lambda)
\asymp_k\lambda_0(bs_0)^k$.

A perturbed first divided difference costs an extra
$O(b/\alpha_*)$. Each Neumann correction costs at most the same.
Replacing $r\ge2$ first divided differences and their intermediate
inverses by a single span-$r$ divided difference costs, relative to
that main scale,
\[
 O\left(\alpha_*^{-1}s_0^{-(r-1)}\right).
\]
If $\lambda_0\ge n^{c/2}/\alpha_*$, every non-main product therefore
has at least one power-small factor, while the other replacement
factors are bounded. Summing the fixed powers of $\log n$ and
charging to the largest block bounds these terms by $o(\mathcal B)$.

Below this threshold, bound absolutely instead. A span-one block,
its following inverse, and one assigned factor $b$ cost
$O(n^{-c/2})$. A higher span $r$ with its following inverse and
$b^r$ costs
\[
 O(b^r\alpha_*^{-1}n^{c/2}/\alpha_*)=O(n^{-rc/2});
\]
the omitted extra factor $\alpha_*^{r-2}n^{-(r-1)c/2}$ is bounded.
Node corrections are bounded as well. The sum of all such tuples is
at most
$n^{1+o(1)}(n^{c/2}/\alpha_*)n^{-kc/2}=o(1)$, by the choice of $k$.
This proves the comparison with the linear pencil.

\emph{Spectral tail and transfer to auxiliaries.}
It follows that $\mathcal B\le C_kn^{.12a}/\alpha_*$. For
$\lambda\ge C/b$, the product in \eqref{eq:original-9} is bounded
below by a positive constant depending only on $k$ and $C$. Hence
the sum above that threshold, and then the norm, are at most
$n^{v_*}/\alpha_*$ for all sufficiently large $n$. Passing between
$S$ and $S_h$ only changes a fixed factor $m_h$. At lower auxiliary
nodes, $M(h-x)^{-1}\preceq M(h)^{-1}$ and
$B_{h-x}^+=(1+o(1))B_h^+$ in relative positive order. The min--max
principle transfers the ordered eigenvalue bounds, increasing the
constant in the tail cutoff.

\emph{Local quadratic estimate.}
The spectral norm bound is now available at all tracked points. We use
it to compare the inverse at a primary with its nearby auxiliary.
If $f$ is a global
coordinate or a centered stub, the path estimates give
$P_{ff}=\|f_\perp\|^2+O(n^{-c_0'})$ and a power-small derivative of
the nonconstant part. Compare $R_h$ and $R_{h-b'}$, where
$b'=n^{-.16a}\alpha$. Their precision difference is
\[
 b'(B_h^+)^{1/2}(I+O(b'/\alpha))(B_h^+)^{1/2}.
\]
The proven norm bound gives
$b'\|S_h\|\le n^{(.125-.16)a}$, a negative power. The resolvent
identity therefore gives the positive quadratic-form comparison
\[
 R_h-R_{h-b'}
 =b'R_h(B_h^+)^{1/2}(I+O(n^{-c'}))(B_h^+)^{1/2}R_h
\]
for some fixed $c'>0$. On a diagonal, the entry error divided by
$b'$ is at most
$Cn^{.53a}/(\alpha_*\sqrt{l_0})=O(n^{-.045a})$.
Using $m_h'/m_h^2=1+h^{-2}$ yields \eqref{eq:original-10} at the
primary.
For a lower auxiliary point at distance $x$, the normalized
precision increment is $x(I+O(x/\alpha))$. Put
$\widetilde S=(B_h^+)^{1/2}R_h(B_h^+)^{1/2}\succeq0$. The map on
the vectors $(B_h^+)^{1/2}R_hf$ has the form
$[I+x\widetilde S T_x]^{-1}$, where
$T_x=I+O(x/\alpha)$ is positive definite. Indeed,
\[
 [I+x\widetilde S T_x]^{-1}
 =T_x^{-1/2}
   [I+xT_x^{1/2}\widetilde S T_x^{1/2}]^{-1}T_x^{1/2}.
\]
The middle factor has norm at most one, so the full norm is at most
$\sqrt{\|T_x\|\|T_x^{-1}\|}=1+O(x/\alpha)$.
Moreover, $B_{h-x}^+$ and $m_{h-x}$ differ from
$B_h^+$ and $m_h$ by relative $O(x/\alpha)$, respectively. Since
$x/\alpha$ is power-small at every auxiliary, these comparisons
transfer \eqref{eq:original-10} to its own normalization at $h-x$,
after reducing $c_1$.
In particular $\sum_xw_i(x)^2\le C$, global unweighted row squares
are at most $C/\alpha$, and $[(\Xi^TW\Xi)^2]_{xx}\le C$.

\emph{Averaged trace identity.}
Take the $\mathcal D$-weighted trace of the same $b'$ resolvent
difference. Its trace-law error after division by $b'l$ is at most
$Cn^{(.09+.16)a}/(l\alpha_*\alpha)$, a negative power. Positivity
and the relative error in the quadratic-form comparison, together
with \eqref{eq:original-10}, give \eqref{eq:original-11}.
Summing the global and stub versions of \eqref{eq:original-10}
also gives $\mathcal T\le C\min(n,l/\alpha)$.

\emph{Moment bounds.}
We now combine the local quadratic estimate with the spectral bulk
and tail bounds. To express the remaining moments in a common matrix,
introduce
$Y=[\Xi,h^{-1}\sqrt\alpha I]$, with the identity projected onto
$\mathbf1^\perp$, and $\mathcal A=Y^TWY$. Its nonzero eigenvalues
are those of $S_h$. Its stub block is $\Gamma$, and the stub block
of $\mathcal A^2$ is $K_{\rm all}+h^{-2}\alpha Q_{\rm all}$,
of trace $O(l)$ by \eqref{eq:original-10}. Split $\mathcal A$ into
its bulk and tail spectral parts. The bulk norm is $O(n^c/\alpha)$;
the tail trace is $O(n^{v_*}/\alpha)$. Their squared parts have
norm $O(n^{2c}/\alpha^2)$ and trace $O(n^{2v_*}/\alpha^2)$,
respectively. Positive compression preserves these bounds.
For a positive bulk compression $C_b$, use
$\operatorname{Tr}C_b^2\le\|C_b\|\operatorname{Tr}C_b$;
for a tail compression $C_t$, use
$\operatorname{Tr}C_t^2\le(\operatorname{Tr}C_t)^2$.
This proves the first estimate in \eqref{eq:original-12}. For the mixed
estimate, split
also the stub block of $\mathcal A$; its bulk contributes
$O(n^cl/\alpha)$, and all terms involving its tail are bounded by
$O(n^{3v_*}/\alpha^3)$, since $c<v_*$. Positive order gives the
last operator bound. Together these estimates prove
\eqref{eq:original-12}.
\end{proof}

\subsection{The estimates passed between pairing stages}

The next two sections establish the hypotheses used above by induction.
The following table records their bounds at the two stopping points;
these are conclusions of Propositions~\ref{prop:early-phase} and
\ref{prop:cleanup}, not additional assumptions. Positivity and the first
three bounds hold for both signs at every tracked primary and auxiliary
parameter. The last bound is needed only at the primaries, and $c_T>0$
denotes a fixed constant.
\begin{center}
\begin{tabular}{@{}lll@{}}
\toprule
Quantity&After early pairing, $l=l_1$&After cleanup, $l=l_0$\\
\midrule
$|L_1|$&$\le n^{.075a}/\alpha_*$&$\le n^{.09a}/\alpha_*$\\
$|L_0|$&$\le n^{.105a}/\alpha_*^2$&$\le n^{.12a}/\alpha_*^2$\\
$\max_{i,j}|E_{ij}|$&$\le n^{.33a}/\sqrt{l_0}$&$\le n^{.37a}/\sqrt{l_0}$\\
$\mathcal T$&$\ge c_T n^{-\theta/2}(n^{-1}+\alpha/l_1)^{-1}$&
 $\ge c_T n^{-\theta}(n^{-1}+\alpha/l_0)^{-1}$\\
\bottomrule
\end{tabular}
\end{center}
The entry labels are global coordinates and surviving centered stubs.
The early bounds leave a positive power of slack inside the caps
\eqref{eq:original-8}; cleanup retains those caps. The structural accounts
also leave fewer than $b$ bad stubs at $l_1$, in the sense of
Lemma~\ref{lemma:structure}. At $l_0$ the unresolved vertices are distinct
and have no nonempty actual nonbacktracking join or return of length at
most $R_0$. Both stages also retain the stub diagonal-square bound
$\sum_{x\in S}E_{xx}^2\le n^\eta$, with a prescribed sufficiently
small fixed loss exponent, and the coefficient and path-multiplicity
bounds. The complete inputs to boundary completion are listed at the
start of Section~\ref{sec:boundary-completion}; only primary parameters
are retained there.
